% GENERATED FILE. Edit canonical lessons, then run npm run generate:books.
% canonical-source-hash: fnv1a64:99780350026f356e
% canonical-lessons: TE-C333-thana, TE-C333-kesu, TE-C333-daryaptu, TE-C333-jada, TE-C333-velimudra, TE-R333-review-326-333

\chapter{At the Police Station}
\label{ch:te-at-the-police-station}
\hlchaptermodality{333}

\begin{chapteropening}
\textbf{By the end of this chapter:} \emph{I can say a police station, a case, an investigation, a trace and a fingerprint.}

Complete the last lesson of chapter 333: I can say a police station, a case, an investigation, a trace and a fingerprint.
\end{chapteropening}

\section[\texorpdfstring{\te{ఠాణా} (ṭhāṇā)}{ఠాణా (ṭhāṇā)}]{\te{ఠాణా} (ṭhāṇā) — a police station}

\label{lesson:TE-C333-thana}



\begin{quote}
Before the new one: say the Telugu for a moment, then the Telugu for to yell.
\end{quote}

\subsection*{You'll want to know: \te{ఠాణా}}
\textbf{\te{ఠాణా}} — \emph{ṭhāṇā} — \textquotedblleft{}a police station\textquotedblright{}.

\begin{quote}
\te{వాళ్ళు} \te{దొంగతనం} \te{గురించి} \te{ఠాణాలో} \te{చెప్పారు}.
\end{quote}

\begin{quote}
\emph{vāḷḷu doṅgatanaṁ guriñci ṭhāṇālō ceppāru.}
\end{quote}

\begin{quote}
They reported the theft at the police station.
\end{quote}

An older word for the \te{పోలీసు} \te{స్టేషన్}; the news still uses it.

\begin{cousinweb}
\textbf{\te{ఠాణా}} (\emph{ṭhāṇā}) is a loan from Hindustani \emph{thānā}, \textquotedblleft{}a post, a station\textquotedblright{}, which goes back to Sanskrit \emph{sthāna}, \textquotedblleft{}a place\textquotedblright{} (CDIAL 13753).
\end{cousinweb}

\begin{grammarlens}[title={What you've built}]
One more word for telling what happened.
\end{grammarlens}

\subsection*{Your turn}
\begin{itemize}
\raggedright
  \item \emph{Say it:} \emph{ṭhāṇā}
  \item \emph{Say it:} \emph{ṭhāṇā}, once more
  \item \emph{Say it:} \emph{vāḷḷu doṅgatanaṁ guriñci ṭhāṇālō ceppāru.} — They reported the theft at the police station.
  \item \emph{Recall:} say the Telugu for a moment, then the Telugu for to yell, then say \emph{ṭhāṇā} again
\end{itemize}

\begin{tcolorbox}[breakable,colback=teal!4,colframe=teal!35!black,title={Before you move on}]
What does \te{ఠాణా} mean? (\textquotedblleft{}A police station\textquotedblright{}.) Say it once more.
\end{tcolorbox}
\section[\texorpdfstring{\te{కేసు} (kēsu)}{కేసు (kēsu)}]{\te{కేసు} (kēsu) — a case (legal)}

\label{lesson:TE-C333-kesu}



\begin{quote}
Before the new one: say the Telugu for to yell, then the Telugu for a police station.
\end{quote}

\subsection*{You'll want to know: \te{కేసు}}
\textbf{\te{కేసు}} — \emph{kēsu} — \textquotedblleft{}a case (legal)\textquotedblright{}.

\begin{quote}
\te{పోలీసులు} \te{దొంగ} \te{మీద} \te{కేసు} \te{పెట్టారు}.
\end{quote}

\begin{quote}
\emph{pōlīsulu doṅga mīda kēsu peṭṭāru.}
\end{quote}

\begin{quote}
The police filed a case against the thief.
\end{quote}

Plural: \te{కేసులు}, \emph{kēsulu}.

\te{కేసు} \te{పెట్టు}, to file a case: literally to put a case on someone.

\begin{cousinweb}
\textbf{\te{కేసు}} (\emph{kēsu}) is a loan from English \emph{case}, with the final -\te{ఉ} that many loans take, as in \te{బస్సు}.
\end{cousinweb}

\begin{grammarlens}[title={What you've built}]
One more word for telling what happened.
\end{grammarlens}

\subsection*{Your turn}
\begin{itemize}
\raggedright
  \item \emph{Say it:} \emph{kēsu}
  \item \emph{Say it:} \emph{kēsu}, once more
  \item \emph{Say it:} \emph{pōlīsulu doṅga mīda kēsu peṭṭāru.} — The police filed a case against the thief.
  \item \emph{Recall:} say the Telugu for to yell, then the Telugu for a police station, then say \emph{kēsu} again
\end{itemize}

\begin{tcolorbox}[breakable,colback=teal!4,colframe=teal!35!black,title={Before you move on}]
What does \te{కేసు} mean? (\textquotedblleft{}A case (legal)\textquotedblright{}.) Say it once more.
\end{tcolorbox}
\section[\texorpdfstring{\te{దర్యాప్తు} (daryāptu)}{దర్యాప్తు (daryāptu)}]{\te{దర్యాప్తు} (daryāptu) — an investigation}

\label{lesson:TE-C333-daryaptu}



\begin{quote}
Before the new one: say the Telugu for a police station, then the Telugu for a case.
\end{quote}

\subsection*{You'll want to know: \te{దర్యాప్తు}}
\textbf{\te{దర్యాప్తు}} — \emph{daryāptu} — \textquotedblleft{}an investigation\textquotedblright{}.

\begin{quote}
\te{పోలీసులు} \te{ఆ} \te{సంఘటనపై} \te{దర్యాప్తు} \te{మొదలుపెట్టారు}.
\end{quote}

\begin{quote}
\emph{pōlīsulu ā saṅghaṭanapai daryāptu modalupeṭṭāru.}
\end{quote}

\begin{quote}
The police began an investigation into the incident.
\end{quote}

\te{దర్యాప్తు} \te{చేయు}, to investigate.

\begin{cousinweb}
\textbf{\te{దర్యాప్తు}} (\emph{daryāptu}) is a loan from Persian \emph{daryāft}, \textquotedblleft{}finding out, inquiry\textquotedblright{}.
\end{cousinweb}

\begin{grammarlens}[title={What you've built}]
One more word for telling what happened.
\end{grammarlens}

\subsection*{Your turn}
\begin{itemize}
\raggedright
  \item \emph{Say it:} \emph{daryāptu}
  \item \emph{Say it:} \emph{daryāptu}, once more
  \item \emph{Say it:} \emph{pōlīsulu ā saṅghaṭanapai daryāptu modalupeṭṭāru.} — The police began an investigation into the incident.
  \item \emph{Recall:} say the Telugu for a police station, then the Telugu for a case, then say \emph{daryāptu} again
\end{itemize}

\begin{tcolorbox}[breakable,colback=teal!4,colframe=teal!35!black,title={Before you move on}]
What does \te{దర్యాప్తు} mean? (\textquotedblleft{}An investigation\textquotedblright{}.) Say it once more.
\end{tcolorbox}
\section[\texorpdfstring{\te{జాడ} (jāḍa)}{జాడ (jāḍa)}]{\te{జాడ} (jāḍa) — a trace; a sign}

\label{lesson:TE-C333-jada}



\begin{quote}
Before the new one: say the Telugu for a case, then the Telugu for an investigation.
\end{quote}

\subsection*{You'll want to know: \te{జాడ}}
\textbf{\te{జాడ}} — \emph{jāḍa} — \textquotedblleft{}a trace; a sign\textquotedblright{}.

\begin{quote}
\te{దొంగ} \te{జాడ} \te{ఎక్కడా} \te{దొరకలేదు}.
\end{quote}

\begin{quote}
\emph{doṅga jāḍa ekkaḍā dorakalēdu.}
\end{quote}

\begin{quote}
No trace of the thief was found anywhere.
\end{quote}

\te{జాడ} \te{తెలియదు}, the trace is not known: nobody knows where someone went.

\begin{cousinweb}
\textbf{\te{జాడ}} (\emph{jāḍa}) is native Dravidian: DEDR 2442 lists Telugu \emph{jāḍa}, \textquotedblleft{}a sign, trace, track\textquotedblright{}, with Tamil \emph{cāṭai}, \textquotedblleft{}a hint, a feature\textquotedblright{}.
\end{cousinweb}

\begin{grammarlens}[title={What you've built}]
One more word for telling what happened.
\end{grammarlens}

\subsection*{Your turn}
\begin{itemize}
\raggedright
  \item \emph{Say it:} \emph{jāḍa}
  \item \emph{Say it:} \emph{jāḍa}, once more
  \item \emph{Say it:} \emph{doṅga jāḍa ekkaḍā dorakalēdu.} — No trace of the thief was found anywhere.
  \item \emph{Recall:} say the Telugu for a case, then the Telugu for an investigation, then say \emph{jāḍa} again
\end{itemize}

\begin{tcolorbox}[breakable,colback=teal!4,colframe=teal!35!black,title={Before you move on}]
What does \te{జాడ} mean? (\textquotedblleft{}A trace; a sign\textquotedblright{}.) Say it once more.
\end{tcolorbox}
\section[\texorpdfstring{\te{వేలిముద్ర} (vēlimudra)}{వేలిముద్ర (vēlimudra)}]{\te{వేలిముద్ర} (vēlimudra) — a fingerprint}

\label{lesson:TE-C333-velimudra}



\begin{quote}
Before the new one: say the Telugu for an investigation, then the Telugu for a trace.
\end{quote}

\subsection*{You'll want to know: \te{వేలిముద్ర}}
\textbf{\te{వేలిముద్ర}} — \emph{vēlimudra} — \textquotedblleft{}a fingerprint\textquotedblright{}.

\begin{quote}
\te{తలుపు} \te{మీద} \te{దొంగ} \te{వేలిముద్రలు} \te{ఉన్నాయి}.
\end{quote}

\begin{quote}
\emph{talupu mīda doṅga vēlimudralu unnāyi.}
\end{quote}

\begin{quote}
The thief's fingerprints are on the door.
\end{quote}

Plural: \te{వేలిముద్రలు}, \emph{vēlimudralu}.

\te{వేలు}, a finger (\te{వేలి}- before another word), with \te{ముద్ర}, a seal or a print.

\begin{cousinweb}
\textbf{\te{వేలిముద్ర}} (\emph{vēlimudra}) is a Telugu compound: native \emph{vēlu}, \textquotedblleft{}a finger\textquotedblright{} (DEDR 5409), and Sanskrit \emph{mudrā}, \textquotedblleft{}a seal, a signet\textquotedblright{} (CDIAL 10203).
\end{cousinweb}

\begin{grammarlens}[title={What you've built}]
One more word for telling what happened.
\end{grammarlens}

\subsection*{Your turn}
\begin{itemize}
\raggedright
  \item \emph{Say it:} \emph{vēlimudra}
  \item \emph{Say it:} \emph{vēlimudra}, once more
  \item \emph{Say it:} \emph{talupu mīda doṅga vēlimudralu unnāyi.} — The thief's fingerprints are on the door.
  \item \emph{Recall:} say the Telugu for an investigation, then the Telugu for a trace, then say \emph{vēlimudra} again
\end{itemize}

\begin{tcolorbox}[breakable,colback=teal!4,colframe=teal!35!black,title={Before you move on}]
What does \te{వేలిముద్ర} mean? (\textquotedblleft{}A fingerprint\textquotedblright{}.) Say it once more.
\end{tcolorbox}
\section[(dialogue)]{Witnesses, kings and rumours}

\label{lesson:TE-R333-review-326-333}



\begin{quote}
Say the Telugu for a witness and a war.
\end{quote}

\subsection*{Your turn}
Say these aloud:

\begin{itemize}
\raggedright
  \item a witness, a war, an incident, to steal, to run away
  \item at one time, a king, a queen, a kingdom, to reign over
  \item at first, afterwards, at long last, the next day, to recall
  \item a theft, to break in, to break open, to make off with, to sense
  \item a reporter, a periodical, a headline, a news story, to publish
  \item a flame, an explosion, a life, to rescue, to run for it
  \item already by then, in the meantime, all of a sudden, a moment, to yell
  \item a police station, a case, an investigation, a trace, a fingerprint
\end{itemize}

\begin{tcolorbox}[breakable,colback=teal!4,colframe=teal!35!black,title={Before you move on}]
A witness? (\textbf{\te{సాక్షి}}, \emph{sākṣi}.) A war? (\textbf{\te{యుద్ధం}}, \emph{yuddhaṁ}.) An incident? (\textbf{\te{సంఘటన}}, \emph{saṅghaṭana}.) To steal? (\textbf{\te{దొంగిలించు}}, \emph{doṅgiliñcu}.) To run away? (\textbf{\te{పారిపోవు}}, \emph{pāripōvu}.) At one time? (\textbf{\te{ఒకప్పుడు}}, \emph{okappuḍu}.) A king? (\textbf{\te{రాజు}}, \emph{rāju}.) A queen? (\textbf{\te{రాణి}}, \emph{rāṇi}.) A kingdom? (\textbf{\te{రాజ్యం}}, \emph{rājyaṁ}.) To reign over? (\textbf{\te{ఏలు}}, \emph{ēlu}.) At first? (\textbf{\te{తొలుత}}, \emph{toluta}.) Afterwards? (\textbf{\te{అనంతరం}}, \emph{anantaraṁ}.) At long last? (\textbf{\te{ఎట్టకేలకు}}, \emph{eṭṭakēlaku}.) The next day? (\textbf{\te{మరుసటి} \te{రోజు}}, \emph{marusaṭi rōju}.) To recall? (\textbf{\te{గుర్తుచేసుకో}}, \emph{gurtucēsukō}.) A theft? (\textbf{\te{దొంగతనం}}, \emph{doṅgatanaṁ}.) To break in? (\textbf{\te{చొరబడు}}, \emph{corabaḍu}.) To break open? (\textbf{\te{బద్దలు} \te{కొట్టు}}, \emph{baddalu koṭṭu}.) To make off with? (\textbf{\te{ఎత్తుకుపోవు}}, \emph{ettukupōvu}.) To sense? (\textbf{\te{పసిగట్టు}}, \emph{pasigaṭṭu}.) A reporter? (\textbf{\te{విలేకరి}}, \emph{vilēkari}.) A periodical? (\textbf{\te{పత్రిక}}, \emph{patrika}.) A headline? (\textbf{\te{శీర్షిక}}, \emph{śīrṣika}.) A news story? (\textbf{\te{కథనం}}, \emph{kathanaṁ}.) To publish? (\textbf{\te{ప్రచురించు}}, \emph{pracuriñcu}.) A flame? (\textbf{\te{మంట}}, \emph{maṇṭa}.) An explosion? (\textbf{\te{పేలుడు}}, \emph{pēluḍu}.) A life? (\textbf{\te{ప్రాణం}}, \emph{prāṇaṁ}.) To rescue? (\textbf{\te{రక్షించు}}, \emph{rakṣiñcu}.) To run for it? (\textbf{\te{పరుగులు} \te{తీయు}}, \emph{parugulu tīyu}.) Already by then? (\textbf{\te{అప్పటికే}}, \emph{appaṭikē}.) In the meantime? (\textbf{\te{అంతలో}}, \emph{antalō}.) All of a sudden? (\textbf{\te{హఠాత్తుగా}}, \emph{haṭhāttugā}.) A moment? (\textbf{\te{క్షణం}}, \emph{kṣaṇaṁ}.) To yell? (\textbf{\te{అరుచు}}, \emph{arucu}.) A police station? (\textbf{\te{ఠాణా}}, \emph{ṭhāṇā}.) A case? (\textbf{\te{కేసు}}, \emph{kēsu}.) An investigation? (\textbf{\te{దర్యాప్తు}}, \emph{daryāptu}.) A trace? (\textbf{\te{జాడ}}, \emph{jāḍa}.) A fingerprint? (\textbf{\te{వేలిముద్ర}}, \emph{vēlimudra}.)
\end{tcolorbox}
